\documentclass{article}

\usepackage{fullpage}
\usepackage{amsmath}
\usepackage{amssymb}
\usepackage{xcolor}

\newcounter{placeholderCount}
\newcommand{\placeholder}[1]{{
    \stepcounter{placeholderCount}
    \color{red} [#1]
}}

\newcounter{solutionCount}
\newcommand{\solution}[1]{
    \stepcounter{solutionCount}
    \paragraph{Solution \thesolutionCount.}
    {
        #1
    }
}

\title{
    Deep Learning (Spring 2026) \\
    Solution for Assignment 2
}
\author{
    Pan Changxun
    \\
    2024011323
}

\begin{document}

\maketitle

\solution{
    True. Generative models can be used to classify data by modeling the joint distribution $p(x,y)$ and applying Bayes' theorem to compute $p(y|x)$.
}

\solution{
True. In maximum likelihood training of an energy-based model, the gradient of the log-likelihood with respect to the parameters 
 depends on the gradient of the log partition function, . This term is exactly equal to the expected gradient of the energy function under the model's distribution: . Therefore, instead of explicitly computing the intractable normalizing factor , we can use sampling methods like MCMC (or Contrastive Divergence) to draw samples from  and approximate this expectation empirically.
}

\solution{
False. Stochastic Gradient MCMC is designed to sample from the posterior distribution $P(\theta|X)$ to capture uncertainty, not to solve an optimization problem. Finding the $\arg \max_{\theta} P(\theta|X)$ is Maximum A Posteriori (MAP) estimation, which seeks a single point estimate and is typically solved using standard optimization algorithms like SGD, not sampling methods.
}

\solution{
    \begin{enumerate}
        \item 
        \textbf{For the noisy image:} Initialize the state of the network with the noisy image. Then, update the neurons iteratively (asynchronously) according to the Hopfield update rule $x_i = \text{sign}(\sum_{j} W_{ij} x_j)$. Since the stored (clean) images are designed to be attractors (local minima of the energy function), the network dynamics will naturally reduce the system's energy, moving the network state away from the noisy initialization and converging towards the nearest clean pattern.
        
        \textbf{For the masked image:} Initialize the state of the network such that the unmasked (observed) pixels are set to their ground truth values, and the masked (unobserved) pixels are initialized randomly or to zero. During the update process, \textbf{clamp} (fix) the states of the observed neurons so they do not change, and only apply the update rule to the unobserved neurons. The fixed correct pixels provide a strong contextual bias, guiding the energy minimization process to correctly reconstruct the missing parts of the corresponding stored pattern.

        \item 
        Suppose the Hebbian learning rule is given by $W = \frac{1}{N}\sum_{p=1}^N y_p y_p^T$, where $y_p \in \{-1, 1\}^N$ are mutually orthogonal patterns. 
        
        Let $Y = [y_1, y_2, \cdots, y_N]$ be an $N \times N$ matrix. Then we can rewrite the weight matrix as $W = \frac{1}{N} Y Y^T$. 
        Since the $N$ patterns are mutually orthogonal, we have $Y^T Y = N I$, where $I$ is the $N \times N$ identity matrix. Because $Y$ is a square matrix, $Y^T Y = N I$ implies that $Y$ is invertible with $Y^{-1} = \frac{1}{N} Y^T$. Consequently, it must also hold that $Y (\frac{1}{N} Y^T) = I$, which means $Y Y^T = N I$. 
        
        Therefore, the weight matrix simplifies to:
        $$W = \frac{1}{N} Y Y^T = I$$
        
        Now, consider an arbitrary network state $x \in \{-1, 1\}^N$. If we update any neuron $x_i$, its new state is determined by the local field:
        $$x_i^{new} = \text{sign}\left(\sum_{j=1}^N W_{ij} x_j\right)$$
        Since $W = I$, $W_{ii} = 1$ and $W_{ij} = 0$ for $i \neq j$. Thus:
        $$x_i^{new} = \text{sign}(W_{ii} x_i) = \text{sign}(x_i)$$
        Given that $x_i \in \{-1, 1\}$, we have $\text{sign}(x_i) = x_i$. 
        
        This demonstrates that for \textit{any} arbitrary configuration $x$, the update rule does not change the state of any neuron ($x^{new} = x$). Therefore, every possible state in the $\{-1, 1\}^N$ space is a stable fixed point (a local minimum of the energy function). Hence, the network successfully memorizes all $2^N$ possible patterns.
    \end{enumerate}
}

\solution{
Given the joint probability 
 where , the marginal probability  can be written as:
Substitute $P(v)$ into the loss function $L(W)$:
\begin{align*}
    L(W) &= -\frac{1}{|P|} \sum_{v \in P} \log \left( \frac{\sum_{h} \exp(y^T W y)}{\sum_{y'} \exp(y'^T W y')} \right) \\
    &= -\frac{1}{|P|} \sum_{v \in P} \left[ \log \left( \sum_{h} \exp(y^T W y) \right) - \log \left( \sum_{y'} \exp(y'^T W y') \right) \right]
\end{align*}

Now, take the gradient of $L(W)$ with respect to $W$. Using the chain rule $\nabla_W \log(f(W)) = \frac{\nabla_W f(W)}{f(W)}$ and the matrix calculus identity $\nabla_W \exp(y^T W y) = y y^T \exp(y^T W y)$:
\begin{align*}
    \nabla_W L(W) &= -\frac{1}{|P|} \sum_{v \in P} \left[ \frac{\sum_{h} \nabla_W \exp(y^T W y)}{\sum_{h} \exp(y^T W y)} - \frac{\sum_{y'} \nabla_W \exp(y'^T W y')}{\sum_{y'} \exp(y'^T W y')} \right] \\
    &= -\frac{1}{|P|} \sum_{v \in P} \left[ \frac{\sum_{h} y y^T \exp(y^T W y)}{\sum_{h} \exp(y^T W y)} - \frac{\sum_{y'} y' y'^T \exp(y'^T W y')}{\sum_{y'} \exp(y'^T W y')} \right]
\end{align*}

Notice that the two terms inside the brackets correspond exactly to expected values:
1. The first term is the expectation of $y y^T$ over the conditional distribution $P(h|v) = \frac{\exp(y^T W y)}{\sum_{h} \exp(y^T W y)}$. Thus, it equals $\mathbb{E}_{h|v}[y y^T]$.
2. The second term is the expectation of $y' y'^T$ over the joint distribution $P(y') = \frac{\exp(y'^T W y')}{\sum_{y'} \exp(y'^T W y')}$. Thus, it equals $\mathbb{E}_{y'}[y' y'^T]$.

Substituting these expectations back into the equation yields the final form:
$$\nabla_W L(W) = -\frac{1}{|P|} \sum_{v \in P} \left( \mathbb{E}_{h|v}[y y^T] - \mathbb{E}_{y'}[y' y'^T] \right)$$
which completes the proof.
}

\solution{
    % Gaussian RBM
    \begin{enumerate}
        \item Since the visible units $v$ are continuous variables ($v \in \mathbb{R}^{N_v}$), the conditional distribution $P(v|h)$ is given by:
        \[ P(v|h) = \frac{\exp(-E(v,h))}{\int \exp(-E(v,h)) dv} \]
        Let's isolate the terms in $-E(v,h)$ that depend on $v$:
        \begin{align*}
        -E(v,h) &= -\frac{1}{2}(v-b)^T(v-b) + v^T Wh \\
        &= -\frac{1}{2}v^T v + v^T b - \frac{1}{2}b^T b + v^T Wh \\
        &= -\frac{1}{2}v^T v + v^T(b + Wh) - \frac{1}{2}b^T b
        \end{align*}
        To recognize this distribution, we can complete the square for $v$. Consider a multivariate Gaussian exponent $-\frac{1}{2}(v-\mu)^T \Sigma^{-1} (v-\mu)$ with $\Sigma = I$ and $\mu = b + Wh$:
        \[ -\frac{1}{2}(v - (b+Wh))^T (v - (b+Wh)) = -\frac{1}{2}v^T v + v^T(b+Wh) - \frac{1}{2}(b+Wh)^T(b+Wh) \]
        Comparing this with our $-E(v,h)$, we can see that:
        \[ -E(v,h) = -\frac{1}{2}(v - (b+Wh))^T (v - (b+Wh)) + \text{const}(h) \]
        Where $\text{const}(h)$ absorbs terms independent of $v$. Therefore, $P(v|h)$ is proportional to the exponential of a squared difference, which means $P(v|h)$ follows a multivariate Gaussian distribution with mean $\mu = b + Wh$ and an identity covariance matrix $\Sigma = I$:
        \[ P(v|h) = \mathcal{N}(v; b + Wh, I) \]
        namely, 
        $$P(v|h) = \frac{1}{(2\pi)^{N_v/2}} \exp\left(-\frac{1}{2}(v - (b+Wh))^T (v - (b+Wh))\right)$$

        \item Let's compute the gradient of $b$ with respect to the log-likelihood for a single data point $v$. Using the general property of RBM gradients:
        \begin{align*}
        \nabla_b \log P(v) &= \nabla_b \log \int_{h} e^{-E(v,h)} - \nabla_b \log Z \\
        &= \mathbb{E}_{h|v}[\nabla_b (-E(v,h))] - \mathbb{E}_{v',h}[\nabla_b (-E(v',h))]
        \end{align*}
        First, compute the derivative of the negative energy with respect to $b$:
        \[ \nabla_b (-E(v,h)) = \nabla_b \left( -\frac{1}{2}(v-b)^T(v-b) + v^T Wh \right) = v - b \]
        Substitute this back into the expectation equation:
        \begin{align*}
        \nabla_b \log P(v) &= \mathbb{E}_{h|v}[v - b] - \mathbb{E}_{v',h}[v' - b]
        \end{align*}
        Since $v$ and $b$ are constants with respect to the conditional distribution $P(h|v)$, the first term simplifies to $v - b$. For the second term, $b$ is a constant:
        \begin{align*}
        \nabla_b \log P(v) &= (v - b) - (\mathbb{E}_{v'}[v'] - b) \\
        &= v - b - \mathbb{E}_{v'}[v'] + b \\
        &= v - \mathbb{E}_{v'}[v']
        \end{align*}
    \end{enumerate}
}

\solution{
    \begin{enumerate}
        \item \textbf{No, $D$ and $F$ are not independent.} 
        In an undirected graphical model, two variables are marginally independent if and only if there is no active path between them. Here, there exists a path $D - A - E - F$. Since none of the intermediate nodes ($A$ or $E$) are observed, the path is not blocked. Therefore, information can flow between $D$ and $F$, making them dependent.

        \item To find the unnormalized conditional distribution $P(B,E|A)$, we marginalize out the unobserved variables $C, D$, and $F$ from the joint distribution. We can omit the global normalization factor $Z$:
        \begin{align*}
            P(B,E|A) &\propto \sum_{C,D,F} P(A,B,C,D,E,F) \\
            &\propto \sum_{C,D,F} f_{AD}(A,D) f_{AC}(A,C) f_{AE}(A,E) f_{BC}(B,C) f_{EF}(E,F) \\
            &= f_{AE}(A,E) \left(\sum_F f_{EF}(E,F)\right) \left(\sum_D f_{AD}(A,D)\right) \left(\sum_C f_{AC}(A,C) f_{BC}(B,C)\right)
        \end{align*}
        
        \textbf{Yes, $B$ and $E$ are independent given $A$.} 
        From the equation above, we can see that the unnormalized conditional probability $P(B,E|A)$ factors into a product of a function of $(A,E)$ and a function of $(A,B)$: 
        \[ P(B,E|A) \propto g_1(A,E) \times g_2(A,B) \times g_3(A) \]
        Since it can be factorized this way, $B \perp\!\!\!\perp E \mid A$. Alternatively, using graph separation, node $A$ blocks the only path between $B$ and $E$ (which is $B-C-A-E$).

        \item We are given that the joint distribution can be modeled as a Boltzmann distribution:
        \[ P(A,B,C,D,E,F) = \frac{1}{Z'} \exp(-E(A,B,C,D,E,F)) \]
        We are also given the factorization of the joint distribution:
        \[ P(A,B,C,D,E,F) = \frac{1}{Z} f_{AC}(A,C) f_{AD}(A,D) f_{AE}(A,E) f_{BC}(B,C) f_{EF}(E,F) \]
        Equating the unnormalized forms, we have:
        \[ \exp(-E(A,B,C,D,E,F)) \propto f_{AC}(A,C) f_{AD}(A,D) f_{AE}(A,E) f_{BC}(B,C) f_{EF}(E,F) \]
        Taking the natural logarithm on both sides:
        \[ -E(A,B,C,D,E,F) = \log f_{AC} + \log f_{AD} + \log f_{AE} + \log f_{BC} + \log f_{EF} + C \]
        By defining the energy functions for each clique as $E_{xy}(x,y) = -\log f_{xy}(x,y)$ (and absorbing any constant), we get:
        \[ E(A,B,C,D,E,F) = E_{AC}(A,C) + E_{AD}(A,D) + E_{AE}(A,E) + E_{BC}(B,C) + E_{EF}(E,F) \]
        This completes the proof.
    \end{enumerate}
}

\solution{
    % Importance Sampling
    We should find the variance minimizer for the importance sampling estimator. 
    \[\mathbb{E}_{x\sim p}[f(x)]\approxeq \frac{1}{N} \sum_{x\sim q} \frac{p(x)}{q(x)}f(x)\]
    The variance of this estimator is:
    \begin{align*}
        \text{Var} &= \mathbb{E}\left[\left(\frac{p(x)}{q(x)}f(x) - \mathbb{E}_{x\sim p}[f(x)]\right)^2\right] \\
        &= \mathbb{E}\left[\left(\frac{p(x)}{q(x)}f(x)\right)^2\right] - \left(\mathbb{E}_{x\sim p}[f(x)]\right)^2 \\
        &= \int \left(\frac{p(x)}{q(x)}f(x)\right)^2 q(x) dx - \left(\int f(x)p(x) dx\right)^2 \\
        &= \int \frac{p(x)^2 f(x)^2}{q(x)} dx - \left(\int f(x)p(x) dx\right)^2
    \end{align*}
    we choose $q(x)$ to minimize the variance. This is equivalent to minimizing the first term since the second term does not depend on $q$:
    \begin{align*}
        \min_q \int \frac{p(x)^2 f(x)^2}{q(x)} dx
    \end{align*}
    Using the method of Lagrange multipliers to enforce the constraint $\int q(x) dx = 1$, we set up the Lagrangian:
    \begin{align*}
        \mathcal{L}(q, \lambda) = \int \frac{p(x)^2 f(x)^2}{q(x)} dx + \lambda \left( \int q(x) dx - 1 \right)
    \end{align*}
    Taking the functional derivative with respect to $q(x)$ and setting it to zero gives:
    \begin{align*}
        \frac{\delta \mathcal{L}}{\delta q(x)} = -\frac{p(x)^2 f(x)^2}{q(x)^2} + \lambda = 0
    \end{align*}
    Solving for $q(x)$, we find:
    \begin{align*} 
        q(x) = \sqrt{\frac{p(x)^2 f(x)^2}{\lambda}} = \frac{p(x)|f(x)|}{\sqrt{\lambda}}
    \end{align*}
    To determine $\lambda$, we use the normalization condition:
    \begin{align*}
        \int q(x) dx = 1 \implies \int \frac{p(x)|f(x)|}{\sqrt{\lambda}} dx = 1 \implies \sqrt{\lambda} = \int p(x)|f(x)| dx
    \end{align*}
    Therefore, the optimal proposal distribution that minimizes the variance of the importance sampling estimator is:
    \begin{align*}
        q^*(x) = \frac{p(x)|f(x)|}{\int p(x)|f(x)| dx}
    \end{align*}
}

\solution{
    \begin{enumerate}
        \item For random-walk using Gaussian noise, the proposal distribution is: 
        \[q(s'|s) = \mathcal{N}(s'; s, \Sigma)\]
        Since Gaussian noise is symmetric, we have $q(s'|s) = q(s|s')$. Therefore, the acceptance probability simplifies to:
        \[\alpha(s \to s') = \min\left(1, \frac{p(s')}{p(s)}\right)\]
        The transition probability from $s$ to $s'$ (for $s \neq s'$) is:
        \[T(s \to s') = q(s'|s) \alpha(s \to s')\]
        Now check the detailed balance condition:
        \begin{align*}
        p(s) T(s \to s') &= p(s) q(s'|s) \alpha(s \to s') \\
        &= p(s) q(s'|s) \min\left(1, \frac{p(s')}{p(s)}\right) \\
        &= q(s'|s) \min(p(s), p(s')) \\
        &= q(s|s') \min(p(s), p(s')) \\
        &= p(s') q(s|s') \min\left(1, \frac{p(s)}{p(s')}\right) \\  
        &= p(s') T(s' \to s)
        \end{align*}
        Thus, the detailed balance condition is satisfied.
        
        The Markov chain is ergodic because it is both irreducible and aperiodic. It is irreducible because the Gaussian noise has infinite support, allowing for transitions between any two states. It is aperiodic because there is a non-zero probability of rejection ($\alpha < 1$), which allows the chain to stay in the current state and prevents deterministic cycles.

        \item For Gibbs sampling:
        In Gibbs sampling, we update one variable at a time by sampling from its conditional distribution given the current values of all other variables. The acceptance rate in the Metropolis-Hastings framework is:
        \[\alpha(s \to s') = \min\left(1, \frac{p(s') q(s' \to s)}{p(s) q(s \to s')}\right)\]
        Split $s$ into $s = (s_i, s_{-i})$ where $s_i$ is the variable being updated and $s_{-i}$ are the other variables. The proposal distribution $q(s \to s')$ is the conditional distribution $p(s_i'|s_{-i})$ multiplied by the probability of choosing index $i$, which is $\frac{1}{N}$. 
        
        Since $p(s) = p(s_i, s_{-i}) = p(s_i|s_{-i})p(s_{-i})$ and $s_{-i}$ remains unchanged ($s'_{-i} = s_{-i}$), we have:
        \begin{align*}
        \alpha(s \to s') &= \min\left(1, \frac{\big[p(s_i'|s_{-i})p(s_{-i})\big] \cdot \big[\frac{1}{N}p(s_i|s_{-i})\big]}{\big[p(s_i|s_{-i})p(s_{-i})\big] \cdot \big[\frac{1}{N}p(s_i'|s_{-i})\big]}\right) \\
        &= \min(1, 1) \\
        &= 1
        \end{align*}
        Thus, the acceptance rate is always $1$ for Gibbs sampling. The Markov chain is ergodic if the conditional distributions are strictly positive, which allows for transitions between all valid states.

    \item \textbf{Proof for Cyclic Gibbs Sampling:(Optional)}
        
        To prove that cyclic Gibbs sampling yields the same stationary distribution $p(s)$, we need to show that $p(s)$ is invariant under one full cycle of updates.
        
        Let $T_i(s \to s')$ denote the transition kernel that updates only the $i$-th coordinate by sampling from the conditional distribution $p(s'_i | s_{-i})$, while keeping all other coordinates unchanged ($s'_{-i} = s_{-i}$).
        
        From the definition of conditional probability, we know that for any coordinate $i$, applying $T_i$ leaves the target distribution $p(s)$ invariant. Mathematically, if $s \sim p(s)$, then after the transition $T_i$, the new state $s'$ still follows $p(s')$:
        \[ \int p(s) T_i(s \to s') ds = p(s') \]
        In operator notation, we can write this invariance as:
        \[ p T_i = p \quad \text{for all } i \in \{1, 2, \dots, N\} \]
        
        In cyclic Gibbs sampling, one complete sweep consists of sequentially applying the transition kernels in a fixed order, typically $T_1, T_2, \dots, T_N$. The combined transition kernel for one full cycle is the composition of these individual kernels:
        \[ T_{cyclic} = T_1 \circ T_2 \circ \dots \circ T_N \]
        
        Now, let's apply this full cyclic transition to the target distribution $p$:
        \begin{align*}
            p T_{cyclic} &= p (T_1 T_2 \dots T_N) \\
            &= (p T_1) T_2 \dots T_N \\
            &= p T_2 \dots T_N \quad \text{(since $p T_1 = p$)} \\
            &= (p T_2) \dots T_N \\
            &= p \dots T_N \quad \text{(since $p T_2 = p$)} \\
            &\ \ \vdots \\
            &= p T_N \\
            &= p \quad \text{(since $p T_N = p$)}
        \end{align*}
        
        Because $p T_{cyclic} = p$, the target distribution $p(s)$ is exactly the invariant (stationary) distribution of the cyclic Gibbs sampling chain. 
        
        Given the problem's assumption that the Markov chain can access all states under the fixed ordering (which guarantees ergodicity), the stationary distribution is unique. Therefore, cyclic Gibbs sampling converges to the same stationary distribution $p(s)$ as the random-order Gibbs sampling.
    \end{enumerate}
}

\solution{
    Used Gemini (AI assistant by Google) in the following capacities:
    \begin{itemize}
        \item \textbf{Solution ideation:} Brainstorming approaches and generating initial ideas for problem-solving strategies, particularly for structuring mathematical proofs and derivations.
        \item \textbf{English grammar and writing:} Polishing and refining the English writing throughout the solutions to improve clarity, coherence, and academic tone.
        \item \textbf{\LaTeX{} formatting:} Optimizing typesetting and layout of mathematical expressions, tables, and document structure for better readability and presentation.
        \item \textbf{Result verification:} Cross-checking the correctness of intermediate computations, final answers, and logical consistency of proofs.
    \end{itemize}
    All final solutions were reviewed and verified by the author.
}

\ifnum\value{placeholderCount}>0
    \vfill
    \begin{center}
        \color{red} \framebox{\textbf{WARNING: PLEASE REMOVE ALL PLACEHOLDERS BEFORE SUBMISSION}}
    \end{center}
\fi

\end{document}